\chapter{Численное решение}

В данной главе представлены результаты вычислительных экспериментов для модели реального делового цикла, реализованных на основе построенного в предыдущей главе численного алгоритма. Основная цель главы --- показать, как построенное глобальное решение (метод стрельбы) описывает динамику модели и чем оно отличается от локального приближения.

В рамках практической части работы выполнены следующие шаги:
\begin{enumerate}
    \item реализована численная схема решения детерминированной RBC-модели с подбором начального потребления \(C_0\) методом стрельбы;
    \item для каждого шага по времени реализовано решение неявного уравнения Эйлера для \(C_{t+1}\) методом Ньютона;
    \item рассчитано стационарное состояние модели и построены переходные траектории из разных начальных условий;
    \item построены фазовые графики в координатах \((K_t, C_t)\) и \((K_t, L_t)\);
    \item построены трехмерные траектории для наглядного анализа выхода экономики к стационарной точке;
    \item выполнено визуальное сравнение локального и глобального решений для DSGE-постановки.
\end{enumerate}

Во всех экспериментах использовались параметры:
\[
\beta = 0.95,\quad
\alpha = 0.33,\quad
\delta = 0.1,\quad
\sigma = 1,\quad
\phi = 0.75,\quad
\psi = 0.35.
\]


\section{Фазовые траектории (K, C) и (K, L)}

На рис.~\ref{fig:rbc_phase_kc} и~\ref{fig:rbc_phase_kl} приведены фазовые траектории для детерминированной модели при различных начальных условиях А - уровня технического прогресса. Для экспериментов были выбраны значения А: $0.95, 1.0, 1.05$, вычислены соответствующие стационарные состояния и построены фазовыве траектории с помощью метода стрельбы слева и справа от стационарного состояния. Первый график показывает совместную динамику капитала и потребления \((K_t, C_t)\), второй --- капитала и труда \((K_t, L_t)\).

\begin{figure}[htbp]
    \centering
    \includegraphics[width=0.9\textwidth]{kc.png}
    \caption{Фазовая траектория детерминированной RBC-модели в координатах \((K_t, C_t)\)}
    \label{fig:rbc_phase_kc}
\end{figure}

\begin{figure}[htbp]
    \centering
    \includegraphics[width=0.9\textwidth]{kl.png}
    \caption{Фазовая траектория детерминированной RBC-модели в координатах \((K_t, L_t)\)}
    \label{fig:rbc_phase_kl}
\end{figure}

\FloatBarrier

Полученные графики подтверждают наличие устойчивой переходной траектории к стационарному состоянию при любых значениях параметра А. Подбор \(C_0\) методом стрельбы обеспечивает сходящуюся динамику без взрывного поведения переменных.

Экспериментально доказаны теорические основы макроэкономики: при росте капитала $K$ потребление $C$ у домохозяйства растёт, а труд $L$ падает. 

\section{Трёхмерная траектория}

Для дополнительной проверки устойчивости построены трехмерные седловые траектории из двух различных начальных условий по капиталу: $K_0 = 10*K^*$ и $K_0 = 0.4*K^*$ соответственно. На рис.~\ref{fig:rbc_3d_first} и~\ref{fig:rbc_3d_second} показаны соответствующие переходные траектории. 
Красными точками на этих двух графиках отмечено стационарное состояние(одинаковое для обоих графиков), зелёными - успешная стартовая точка, из которой метод стрельбы сошёлся к стационарному состоянию. Слегка прозрачными линиями - неудачные траектории, которые не сошлись к стационарному состоянию. 

\begin{figure}[htbp]
    \centering
    \includegraphics[width=0.9\textwidth]{rbc_3d.png}
    \caption{Трехмерная седловая траектория при $K_0 = 10*K^*$}
    \label{fig:rbc_3d_first}
\end{figure}

\begin{figure}[htbp]
    \centering
    \includegraphics[width=0.9\textwidth]{rbc_3d_2.png}
    \caption{Трехмерная седловая траектория при $K_0 = 0.4*K^*$}
    \label{fig:rbc_3d_second}
\end{figure}

\FloatBarrier

Обе траектории сходятся к одной и той же стационарной точке, что является ожидаемым результатом для седлово-устойчивой динамической системы в окрестности равновесия. На графиках виден разброс, которые может возникнуть при неудачном выборе начальное потребления $C_0$

\section{Сравнение локального и глобального решения DSGE}

В DSGE-модели решением является распределение траекторий макропеременных, а не одна детерминированная траектория. На рис.~\ref{fig:local_global_dsge} и~\ref{fig:local_global_dsge_2} приведено сравнение локального и глобального подходов по всем десяти эндогенным переменным модели. Для удобства восприятия они разбиты на два рисунка: реальные величины и труд (рис.~\ref{fig:local_global_dsge}) и цены и факторы (рис.~\ref{fig:local_global_dsge_2}). По оси Y указаны проценты отклонения от стационарного состояния (за исключением прибыли \(\mathit{Pr}\), для которой используется отклонение в уровнях), по оси X --- периоды симуляции.

\begin{figure}[htbp]
    \centering
    \includegraphics[width=\textwidth]{comparison_dynare_vs_chebyshev_part1.png}
    \caption{Сравнение локального  и глобального решений на шок \(\mathit{tfp}\) с \(\sigma_{\text{tfp}} = 0{,}2\): реальные величины и труд --- выпуск \(Y\), капитал \(K\), труд \(L\), потребление \(C\), инвестиции \(I\), зарплата \(W\).}
    \label{fig:local_global_dsge}
\end{figure}

\begin{figure}[htbp]
    \centering
    \includegraphics[width=\textwidth]{comparison_dynare_vs_chebyshev_part2.png}
    \caption{Сравнение локального и глобального решений: цены и факторы --- процентная ставка \(R\), прибыль \(\mathit{Pr}\), предельные издержки \(\mathit{MC}\) и совокупная факторная производительность \(\mathit{tfp}\). Для \(\mathit{Pr}\) показано отклонение в уровнях (стационарное значение равно нулю при функции Кобба-Дугласа), для остальных --- проценты от стационарного состояния.}
    \label{fig:local_global_dsge_2}
\end{figure}

\FloatBarrier

Локальное решение, получаемое через линеаризацию вокруг стационарного состояния, хорошо работает в малой окрестности равновесия и отличается вычислительной простотой. Глобальное решение учитывает нелинейность модели на всем диапазоне состояний и более корректно описывает динамику при больших отклонениях от стационарной точки.

Практический вывод состоит в следующем: для быстрого анализа малых шоков локальные методы удобны, но для задач, связанных с существенными отклонениями и исследованием полной переходной динамики, предпочтителен глобальный подход, который вычислительно дорогой и требует больше времени. 

\paragraph{Монетарный шок: сравнение локального и глобального решений.}

Сравнение проводится по импульсным функциям отклика всех 19 эндогенных переменных на единичный монетарный шок. Для удобства восприятия графики разбиты на два графика - реальные величины и труд (рис.~\ref{fig:mon_compare_real}) и цены, прибыль и факторы (рис.~\ref{fig:mon_compare_prices}).

\begin{figure}[htbp]
    \centering
    \includegraphics[width=\textwidth]{compare_orders_mon_part1.png}
    \caption{Монетарный шок, сравнение локального и глобального решений: реальные величины --- выпуск \(Y\), капитал \(K\), труд \(L\), потребление \(C\), инвестиции \(I\), зарплаты \(W\), \(W^{real}\), аренда капитала \(R\).}
    \label{fig:mon_compare_real}
\end{figure}

\begin{figure}[p]
    \centering
    \includegraphics[width=\textwidth, height=0.92\textheight, keepaspectratio]{compare_orders_mon_part2.png}
    \caption{Монетарный шок, сравнение локального и глобального решений: реальная аренда \(R^{real}\), номинальная ставка \(R^b\), реальная \(R^{b,real}\), предельные издержки \(MC\), \(MC^{real}\), прибыль \(\mathit{Pr}\), \(\mathit{Pr}^{real}\), уровень цен \(P\), инфляция \(\pi\), TFP, монетарный мультипликатор \(\mathit{mon}\).}
    \label{fig:mon_compare_prices}
\end{figure}

\FloatBarrier

Анализ результатов показывает, что для большинства <<номинальных>> переменных (\(W\), \(W^{real}\), \(R\), \(R^{real}\), \(MC\), \(MC^{real}\), \(\mathit{Pr}\), \(\mathit{Pr}^{real}\)) линейное и нелинейное приближения дают существенно различающиеся импульсные отклики. В частности, линейное решение предсказывает крайне резкую начальную реакцию (\(W\) опускается до \(-35\%\), \(R\) --- до \(-80\%\), прибыль вырастает до \(+200\%\)), тогда как нелинейное приближение демонстрирует значительно более умеренные отклики того же знака. Это связано с тем, что параметр \(\mu = 30\) обеспечивает сильную нелинейность в уравнении Ротемберга, которую линеаризация 1-го порядка адекватно не передаёт.

Для <<реальных>> величин (\(Y\), \(K\), \(L\), \(C\), \(I\)) обе схемы дают качественно близкие отклики --- сокращение выпуска и труда, плавное восстановление капитала и потребления, --- но и здесь нелинейное приближение даёт количественно меньшие амплитуды.

Полученные результаты подтверждают общий вывод работы: при анализе реакции экономики на крупные монетарные шоки следует использовать решения более высокого порядка, поскольку линейное приближение систематически переоценивает амплитуду откликов <<номинальных>> переменных.
